\section{Proofs}\label{supp:proofs}

This section proves the statements of Sections~\ref{sec:arrow} and~\ref{sec:theory} of the main text, in the order in which they appear, and adds the worked examples they refer to. None of these results is a convergence or generalization theorem for the trained network. Section~\ref{sec:theory} states several results only in words. Three of them receive exact statements as propositions below, and the subsections on the alphabet, the angle and the batch update prove the others. Each subsection is headed with the part of Section~\ref{sec:theory} it supports. Positions are counted from $1$, and the implementation counts from $0$. Adding one to every position changes no ordering, no footrule distance and no aggregate in exact arithmetic, because position sums shift by a common constant.

\subsection{Details for Section~\ref*{sec:arrow}: The Update Rule and Arrow's Axioms}\label{supp:iia}

Every update of a ranking filter merges several rankings into one, so it can be studied with Arrow's impossibility theorem. The weighted Borda count is Pareto efficient and non-dictatorial under the weight condition of Proposition~\ref{prop:borda-iia}, so by Arrow's theorem it violates independence of irrelevant alternatives (IIA). This subsection proves the proposition. It also gives witnesses, small examples on which the violation can be checked by hand. Last, it treats the update as training runs it, with the filter's current order held fixed.

A profile lists one ranking per voter, and two profiles are compared voter by voter. The three axioms are those of Section~\ref{sec:arrow}. Pareto efficiency asks that the collective ranking place $a$ before $b$ whenever every voter does. Non-dictatorship asks that no single voter's ranking be the collective ranking on every profile. IIA asks that the collective order of $a$ and $b$ depend only on how each voter orders $a$ and $b$, so that moving a third item on the ballots cannot change it.

The rule under study is the weighted Borda count~\eqref{eq:borda-update}. With fixed positive weights $\omega_i$, item $v$ receives the score $P(v)=\sum_i\omega_i\operatorname{pos}_{\pi_i}(v)$ over the voters $i$. The collective ranking sorts the items by $P$, with ties broken by the fixed item order. When a batch update is read as an election (Section~\ref{sec:arrow}), the prior $r_j$ has weight one and voter $t$ has weight $\omega_t=|a_t|$.

\begin{proof}[Proof of Proposition~\ref{prop:borda-iia}]
\emph{Universal domain.} $P$ is defined for every profile of complete rankings, so the rule returns a ranking on every profile.

\emph{Pareto.} If every voter places $a$ before $b$, then $\operatorname{pos}_{\pi_i}(a)<\operatorname{pos}_{\pi_i}(b)$ for every $i$. Multiplying by $\omega_i>0$ and summing gives $P(a)<P(b)$, so $a$ precedes $b$ whatever the tie rule.

\emph{Non-dictatorship.} Fix a voter $i$ and let $\Omega_{-i}$ be the combined weight of the other voters. By assumption, $\omega_i\le(V-1)\Omega_{-i}$. Take two items $a$ and $b$, with $a$ the one of higher identifier (ID). Consider the profile in which voter $i$ places $a$ first and $b$ second, while every other voter places $b$ first and $a$ last. Then
\[
P(b)-P(a)=\omega_i(2-1)+\sum_{i'\neq i}\omega_{i'}(1-V)=\omega_i-(V-1)\Omega_{-i}\le0.
\]
So either $b$ precedes $a$ strictly, or the two tie. In a tie, the tie rule puts the lower ID first, and that is $b$. In both cases the collective ranking disagrees with voter $i$, so no voter is a dictator.

\emph{IIA.} Arrow's theorem \citep{arrow1963social} concerns rules on profiles of rankings of at least three alternatives. It states that such a rule with universal domain that satisfies Pareto and is non-dictatorial cannot satisfy IIA. The witnesses below show the violation directly.
\end{proof}

The weight condition is exact. A voter that weighs more than $V-1$ times all others combined is a dictator. Whenever it places $a$ before $b$, $P(b)-P(a)\ge\omega_i-(V-1)\Omega_{-i}>0$. For the prior, whose weight is one, this happens when the total vote weight of the batch is below $1/(V-1)$. Such a batch cannot move the filter. How often a batch reaches this weight in training is reported only for the first of two hidden layers under the corrected relay. There, 58.4 percent of the filter-batches with a nonzero vote reached it on average, and 97.7 percent when these votes were multiplied by 8 (Table~\ref{tab:s-relay-movement}).

\paragraph{Witness without the prior.}
Take three items $a,b,c$ and two voters of weight one. A ballot $abc$ places $a$ first, $b$ second and $c$ third. Compare the profiles
\[
\Pi=(abc,\;bca)\qquad\text{and}\qquad\Pi'=(acb,\;bac).
\]
In both, voter 1 places $a$ before $b$ and voter 2 places $b$ before $a$. Only the position of $c$ relative to $a$ and $b$ changes. The position sums $(P(a),P(b),P(c))$ are $(4,3,5)$ for $\Pi$, giving $bac$, and $(3,4,5)$ for $\Pi'$, giving $abc$. The order of $a$ and $b$ reverses strictly, so the tie rule plays no part. Moving $c$ alone thus reverses the collective order of $a$ and $b$, which IIA forbids.

\paragraph{Witness with the prior.}
The same happens with a prior. Take the prior $r=abc$ with weight one and six votes of weight $\tfrac12$, and compare the profiles
\[
\widetilde\Pi=(abc,abc,abc,bca,bca,bca)\qquad\text{and}\qquad\widetilde\Pi'=(acb,acb,acb,bac,bac,bac).
\]
Every corresponding voter keeps its order of $a$ and $b$, and the prior and weights are identical. The position sums are $(7,\tfrac{13}{2},\tfrac{21}{2})$ for $\widetilde\Pi$, giving $bac$, and $(\tfrac{11}{2},8,\tfrac{21}{2})$ for $\widetilde\Pi'$, giving $abc$. Dividing by the total weight $4$ gives the mean voted positions of~\eqref{eq:rowmean-borda}. The weight condition holds, since $1<2\cdot3$ for the prior and $\tfrac12<2\cdot\tfrac72$ for each vote.

Both witnesses concern the aggregation rule with fixed voters and weights. In training, the accepted targets and their weights depend on the current predictions and on representations that change from batch to batch. And no vote from an earlier batch is retained. The witnesses therefore show what the rule does on such profiles. They do not show that such profiles arise in training, and they imply nothing about classification accuracy.

\paragraph{The domain of the proposition.} The proposition varies every ballot, the prior included. One update, however, holds the prior fixed. The map from the remaining ballots is then defined on every profile of those ballots, so by Arrow's theorem it cannot satisfy all three axioms either. Which axioms it gives up depends on the total vote weight $\Omega$ of the batch, and there are three cases (Figure~\ref{fig:s-weight-regimes}). If $\Omega<1/(V-1)$, the current order decides alone and the map is constant. The constant map satisfies IIA, and no ballot dictates, but it is not Pareto efficient.

If $1/(V-1)<\Omega<V-1$, the map is not Pareto efficient. Take items $a$ and $b$ at the two ends of the prior, and let every ballot place $b$ immediately before $a$. Then $P(b)-P(a)=(V-1)-\Omega>0$. Arrow's theorem, which assumes Pareto efficiency, then says nothing about IIA. But IIA fails directly, so here the map gives up both Pareto efficiency and IIA. Take two items with $a$ before $b$ in the prior, $g$ positions apart, where $\Omega<g<(V-1)\Omega$. For example, $g=1$ works if $\Omega<1$, and $g=\lfloor\Omega\rfloor+1$ otherwise. If every ballot places $b$ immediately before $a$, then $P(b)-P(a)=g-\Omega>0$ and the filter keeps $a$ first. If every ballot places $b$ first and $a$ last, then $P(b)-P(a)=g-(V-1)\Omega<0$ and the filter places $b$ first. Yet no ballot changed its order of $a$ and $b$.

At $V=3$ with $\Omega=1$ this construction fails, because no such $g$ exists, but IIA still fails through the tie rule. Take the prior $abc$, which is also the identifier order, and one ballot of weight one. The ballots $cba$ and $cab$ both place $c$ before $b$. The first gives the position sums $(4,4,4)$, and the tie rule keeps $b$ before $c$. The second gives $(3,5,4)$ and places $c$ before $b$. For every other identifier order, one of the three pairs of the prior gives a witness of the same kind.

If $\Omega>V-1$, the map is Pareto efficient. Unless one ballot outweighs $V-1$ times all other weight, the prior's included, it violates IIA by Arrow's theorem, as the six-vote witness above does with its prior fixed.

At each boundary between these cases the tie rule decides. For example, take $V=3$ and the prior $abc$ of weight one, which is also the identifier order. Add one vote of weight $\tfrac12$, so that $\Omega=1/(V-1)$. Every pair of items keeps a score gap of at least $1-2\cdot\tfrac12=0$ in the order of the prior, and an equality goes to the identifier order. So each of the six possible ballots returns $abc$. This map is constant. It satisfies IIA, and no ballot dictates. But it is not Pareto efficient, since the ballot $cba$ places $c$ before $a$ and the output does not. With a prior that disagrees with the identifier order, the update can move the filter at this boundary. It can do so only through a tie, because every pair keeps a score gap of at least $1-(V-1)\Omega=0$ in the order of the prior. Such a prior has two neighbors, $a$ just before $b$, with $b$ of lower identifier. If every ballot places $b$ first and $a$ last, then $P(b)-P(a)=1-(V-1)\Omega=0$, and the tie rule puts $b$ first. If every ballot places $b$ immediately before $a$, then $P(b)-P(a)=1-\Omega>0$ and $a$ stays first. So IIA fails at this boundary too.

In ArrowFlow's layers each vote weighs less than $2\eta\le0.4$ and a batch has $32$ examples, so $\Omega<12.8<15\le V-1$. No update is therefore Pareto efficient, and an update that can move its filter violates IIA by the constructions above.

% Figure: the three weight ranges of one update with the prior fixed (supplement, Section S1.1). Drawn by hand; it
% restates the cases of the paragraph "The domain of the proposition" (not to scale).
\begin{figure}[!htbp]
\centering
\begin{tikzpicture}[
    >=Stealth,
    lbl/.style={font=\footnotesize, align=center},
    tick/.style={font=\footnotesize},
]
\fill[green!10] (0,0) rectangle (3,0.5);
\fill[orange!15] (3,0) rectangle (10,0.5);
\fill[blue!10] (10,0) rectangle (14,0.5);
\draw[->, thick] (0,0) -- (14.3,0) node[right, font=\footnotesize] {$\Omega$};
\draw[thick] (3,-0.12) -- (3,0.5) node[above, tick] {$1/(V-1)$};
\draw[thick] (10,-0.12) -- (10,0.5) node[above, tick] {$V-1$};
\node[lbl, anchor=south] at (1.5,1.1) {gives up only\\Pareto efficiency,\\so the filter\\never moves};
\node[lbl, anchor=south] at (6.5,1.1) {gives up Pareto efficiency\\and independence, which\\fails directly (at $V=3$,\\$\Omega=1$ through the tie rule)};
\node[lbl, anchor=south] at (12,1.1) {gives up independence:\\Pareto efficient, so\\Arrow's theorem\\forces the failure};
\draw[decorate, decoration={brace, mirror, amplitude=5pt}, thick] (0,-0.35) -- (8.6,-0.35)
      node[midway, below=7pt, lbl] {ArrowFlow's layers: $\Omega<12.8<15\le V-1$};
\end{tikzpicture}
\caption{\textbf{What one update gives up when the current order is held fixed.} One batch update of a filter over $V$ items depends on the total vote weight $\Omega$ of the batch, drawn here on an axis that is not to scale. Below $1/(V-1)$ the current order decides alone, so the update never moves the filter. This constant map satisfies independence of irrelevant alternatives, and no ballot dictates, but it is not Pareto efficient. Between $1/(V-1)$ and $V-1$ the update gives up both Pareto efficiency and independence. Independence fails directly by the construction in this subsection, and at $V=3$ with $\Omega=1$ through the tie rule. Figure~\ref{fig:iia} of the main text shows a case from this middle range. Above $V-1$ the update is Pareto efficient, and Arrow's theorem forces the failure unless one ballot outweighs $V-1$ times all other weight, the prior's included. In ArrowFlow's layers every update lies in the first two ranges. At each boundary the tie rule decides.}
\label{fig:s-weight-regimes}
\end{figure}


\paragraph{The forward pass.} Section~\ref{sec:arrow} observes that the forward pass orders any two filters of a layer independently of the other filters. This is not Arrow's independence axiom, because the forward pass aggregates no profile of rankings.

\subsection{Score-Gap Stability}

This subsection proves Proposition~\ref{prop:score-stability}. If every score moves by less than half the smallest gap between two scores, the encoded ranking stays the same. The proof is short, because the score-gap condition bounds the change in every pairwise difference of the scores at once.

\begin{proof}[Proof of Proposition~\ref{prop:score-stability}]
Take any pair with $z_i<z_j$. Then
\[
(z_j+\varepsilon_j)-(z_i+\varepsilon_i)\ge z_j-z_i-2\|\varepsilon\|_\infty\ge\delta_{\min}(z)-2\|\varepsilon\|_\infty>0 .
\]
Every strict pairwise order is preserved, so the sorted order of the coordinates, and with it $\argsort$, is unchanged. The strict inequality in the hypothesis rules out a perturbation that creates a tie.
\end{proof}

\subsection{Details for Section~\ref*{sec:theory-stability}: Gaussian Perturbation through an Affine Encoder}\label{supp:gaussian}

Under Gaussian noise the score-gap condition cannot hold for sure, because Gaussian noise exceeds any bound with positive probability. This subsection therefore bounds the probability that Gaussian noise on the input changes the encoded ranking. Section~\ref{sec:theory-stability} summarizes this result in words. For an encoder of degree one the scores are affine in the input, that is, a linear function of it plus a constant. So Gaussian input noise stays Gaussian at the score level.

\begin{proposition}[Gaussian perturbation through an affine encoder]\label{prop:gaussian-stability}
Let $z_i(x)=w_i^{\top}x+b_i$ for $i=1,\dots,e$, with any fitted standardization folded into $w_i$ and $b_i$. Assume that $z(x)$ has distinct coordinates. Let $\varepsilon\sim\mathcal N(0,\Sigma)$ with $\Sigma$ positive semidefinite, and put $\delta_{ij}=|z_i(x)-z_j(x)|$ and $v_{ij}=(w_i-w_j)^{\top}\Sigma\,(w_i-w_j)$. Then
\begin{equation}\label{eq:gaussian-pair-bound}
\Pr\bigl[\argsort(z(x+\varepsilon))\neq\argsort(z(x))\bigr]
\le\min\Bigl\{1,\;\sum_{i<j,\;v_{ij}>0}\Phi_{\mathcal N}\bigl(-\delta_{ij}/\sqrt{v_{ij}}\bigr)\Bigr\},
\end{equation}
where $\Phi_{\mathcal N}$ is the standard normal distribution function, and a pair with $v_{ij}=0$ contributes zero.
\end{proposition}

The bound uses each pair's own gap and variance, and it allows correlated input noise. Each pair of scores reverses with a normal tail probability, and the union bound combines the pairs. The union bound says that the probability that at least one of several events occurs is at most the sum of their probabilities. The last paragraph of this subsection shows that the sum is also the expected Kendall distance between the clean and perturbed rankings. The Kendall distance counts the pairs that two rankings order differently.

\begin{proof}[Proof of Proposition~\ref{prop:gaussian-stability}]
Since $z_i(x+\varepsilon)=z_i(x)+w_i^{\top}\varepsilon$, a pair with $z_i(x)<z_j(x)$ has the perturbed difference $\delta_{ij}+(w_j-w_i)^{\top}\varepsilon$. Its random part is a normal variable with mean zero and variance $v_{ij}$. If $\argsort$ changes, some pair that is strictly ordered at $x$ is reversed or tied at $x+\varepsilon$. For a pair with $v_{ij}>0$ the reversal has probability $\Phi_{\mathcal N}(-\delta_{ij}/\sqrt{v_{ij}})$ and the tie has probability zero. For a pair with $v_{ij}=0$ the random part vanishes almost surely, that is, with probability one, so neither can happen. The union bound over the pairs gives the sum in~\eqref{eq:gaussian-pair-bound}, and a probability is at most one. The union bound needs no independence between pairs.
\end{proof}

For isotropic noise, which has the same variance in every direction, $\Sigma=\sigma^2I$ gives $v_{ij}=\sigma^2\|w_i-w_j\|_2^2$. Treating the projected coordinates as independent would give a different variance and is unnecessary.

\paragraph{The sum as an expected distance.} Let $S$ be the sum of pair probabilities in~\eqref{eq:gaussian-pair-bound}, before the cap at one. Each of its terms is exactly the probability that its pair reverses. A reversed pair is a pair that the clean and perturbed rankings order differently. By linearity of expectation, $S$ is therefore the expected Kendall distance $\E[d_K]$ between the two rankings, and $d_F\le2d_K$ \citep{diaconis1977footrule} gives $\E[d_F]\le2S$. This extends the bound to the first ranking layer. Suppose the responses of a first ranking layer at $\pi_x$, the footrule distances from $\pi_x$ to the layer's filters, are distinct, with smallest consecutive gap $g_{\min}$. By Proposition~\ref{prop:layer-margin}(ii) and Markov's inequality, the layer keeps its output ordering at $x+\varepsilon$ except with probability at most $4S/g_{\min}$. Markov's inequality bounds the probability that a nonnegative variable is at least some level by its mean divided by that level. A changed output also needs a changed encoding, so the probability that the output changes is at most $\min\{1,S,4S/g_{\min}\}$. At the layer sizes used here, however, the responses are rarely all distinct (Section~\ref{supp:margin}). Without that condition the bound is $\min\{1,S\}$, the bound for the encoding itself.

\subsection{Local Sensitivity of Polynomial Encoders}\label{supp:poly}

The bound of Section~\ref{supp:gaussian} needs an encoder of degree one, whose scores are affine in the input. An encoder of higher degree makes each score a polynomial in the input. This subsection shows that for small noise the variance of a change in a score difference is still set, to leading order, by the gradient of that difference.

Let $f\colon\R^d\to\R^e$ be a fitted encoder of degree $p$, the composition of polynomial expansion, standardization and projection, so that each score $f_i$ is a polynomial in $x$. Fix $x$ and a positive semidefinite matrix $\Sigma_0$, and let $\varepsilon=\sigma\zeta$ with $\zeta\sim\mathcal N(0,\Sigma_0)$. For the score difference $\Delta_{ij}(x)=f_i(x)-f_j(x)$,
\begin{equation}\label{eq:poly-local}
\operatorname{Var}\bigl[\Delta_{ij}(x+\varepsilon)-\Delta_{ij}(x)\bigr]=\sigma^2\,\nabla\Delta_{ij}(x)^{\top}\Sigma_0\,\nabla\Delta_{ij}(x)+O(\sigma^4)\qquad(\sigma\to0),
\end{equation}
where the constant in the remainder may depend on $x$, $f$ and $\Sigma_0$.

\begin{proof}
The polynomial $\Delta_{ij}$ has a finite Taylor expansion,
\[
\Delta_{ij}(x+\sigma\zeta)-\Delta_{ij}(x)=\sigma\,\nabla\Delta_{ij}(x)^{\top}\zeta+\sigma^2 q_2(\zeta)+\sigma^3 q_3(\zeta)+\dots,
\]
with $q_2$ homogeneous quadratic, $q_3$ homogeneous cubic and finitely many further terms of higher degree in $\zeta$. A homogeneous polynomial is one whose terms all have the same total degree. The variance of the linear term is the quadratic form in~\eqref{eq:poly-local}. Its covariance with $q_2(\zeta)$ vanishes. Because $q_2$ is homogeneous of degree two, its product with the linear term is homogeneous of degree three. And every moment of odd total degree of a Gaussian vector with mean zero is zero. Every remaining contribution to the variance is of order $\sigma^4$ or higher in $\sigma$, and all Gaussian moments are finite.
\end{proof}

The gradient in~\eqref{eq:poly-local} is built from the derivatives of monomials. For a monomial $x^{\alpha}=\prod_i x_i^{\alpha_i}$, a product of powers of the features, the derivative with respect to $x_i$ is
\[
\frac{\partial x^{\alpha}}{\partial x_i}=\alpha_i\,x_i^{\alpha_i-1}\prod_{i'\neq i}x_{i'}^{\alpha_{i'}}
\]
when $\alpha_i>0$ and zero otherwise. The expression covers repeated powers and does not divide by $x_i$, so it is defined at zero. For example, $x_i^{p}$ has derivative $p\,x_i^{p-1}$. The gradient of a score difference combines these derivatives with the fitted standardization and projection coefficients. So the sensitivity of a single monomial before projection is not enough to describe that of a score.

Equation~\eqref{eq:poly-local} is a local variance expansion, valid for small noise. A nonlinear polynomial of Gaussian noise is generally not Gaussian and may have a nonzero mean shift. So inserting this variance into~\eqref{eq:gaussian-pair-bound} would not give an exact bound for noise of finite size.

\subsection{Details for Section~\ref*{sec:theory-alphabet}: The Ordinal Alphabet and Permutation Cones}\label{supp:alphabet}

This subsection proves the statements of Section~\ref{sec:theory-alphabet} about the alphabet, the set of rankings that the encoder can emit. Ranking keeps order and discards magnitude. Applying one strictly increasing function to every coordinate leaves $\argsort(z)$ unchanged. This holds for the vector that is sorted. It does not hold for raw features transformed before expansion and projection. For scores without ties, $\argsort(z)$ takes one of $e!$ values, so the entropy of the encoded ranking obeys
\begin{equation}\label{eq:alphabet}
H(\pi_x)\le\log_2(e!),
\end{equation}
with equality only when all $e!$ orderings are equally likely.

All score vectors without ties that share one ranking form a cone (Figure~\ref{fig:cones} of the main text). The encoding is constant inside a cone and jumps at its walls. Each ranking layer sorts distances, so it is piecewise constant too. A fixed projection may land in only some cones, and at degree one their number can be counted (the last paragraph of this subsection). So~\eqref{eq:alphabet} bounds the alphabet and not the information the data carry, and it predicts neither classification error nor network capacity. At fixed $e$, polynomial expansion leaves the bound $e!$ unchanged, but it can enlarge the set of rankings that the encoder reaches. And the preimages of the cones, the sets of inputs that land in each cone, need not be convex. Training the layers that follow a fitted encoder does not move the cone walls. The rest of this subsection proves these statements.

Two score vectors without ties have the same $\argsort$ if and only if $z_i-z_j$ has the same sign in both for every pair $i\neq j$. For a permutation $\pi$ of the $e$ coordinates, the set $\mathcal{C}_\pi=\{z: z_{\pi[1]}<z_{\pi[2]}<\dots<z_{\pi[e]}\}$ is defined by $e-1$ strict linear inequalities, so it is open and convex. It is closed under positive scaling, which makes it a cone. It is also closed under adding a multiple of the all-ones vector. The $e!$ sets $\mathcal{C}_\pi$ are pairwise disjoint. The union of the hyperplanes $\{z_i=z_j\}$ has Lebesgue measure zero, that is, zero volume, and its complement is the union of the cones. Since $\argsort$ is constant on each $\mathcal{C}_\pi$ and takes different values on different cones, the encoded ranking takes at most $e!$ values. The entropy of a random variable with at most $e!$ values is at most $\log_2(e!)$, with equality exactly when all $e!$ values are equally likely. This is~\eqref{eq:alphabet}. It remains to show that a strictly increasing function leaves $\argsort(z)$ unchanged. Such a function preserves every strict inequality and every equality between coordinates. So all pairwise comparisons are unchanged, and therefore so are the sorted order and the tie rule's decisions.

\paragraph{The rankings a view reaches at degree one.} At degree one the scores of every view are linear in the standardized features $q\in\R^d$, with no offset (Section~\ref{supp:encoder}). So each score is $z_i=q^{\top}w_i$ for a fitted vector $w_i\in\R^d$. The encoded ranking is then the order in which the linear function $w\mapsto q^{\top}w$ sorts the $e$ points $w_1,\dots,w_e$ of $\R^d$. \citet{cover1967orderings} counted such orders. For $e$ points in general position in $\R^d$, linear functions induce exactly $Q(e,d)$ orders, where $Q(1,d)=1$, $Q(e,1)=2$ for $e\ge2$, and $Q(e+1,d)=Q(e,d)+e\,Q(e,d-1)$ for $d\ge2$. Points in other positions give at most as many. The columns of a Gaussian matrix are in general position with probability one, so a random view reaches all $Q(e,d)$ orders as $q$ ranges over $\R^d$. For $e=16$ and $d=4$ that is $437{,}040$ rankings, against $16!\approx2.1\times10^{13}$. A polynomial expansion makes the scores linear in more coordinates, which raises this bound.

\subsection{Angle Property of the Random Projection}\label{supp:angle}

For the random projection strategy, the expected disagreement between the rankings of two inputs is set by the angle between their expanded, standardized feature vectors. So inputs that point in similar directions get similar rankings on average. This subsection states and proves the exact relation.

Let $q,q'\in\R^{d'}$ be the two feature vectors, both nonzero, with angle $\theta\in[0,\pi]$. Let $W\in\R^{d'\times e}$, $e\ge2$, have independent standard normal entries. Put $\kappa=\argsort(qW)$ and $\kappa'=\argsort(q'W)$. Let $d_K(\kappa,\kappa')$ be the Kendall distance \citep{kendall1938}, the number of coordinate pairs that the two rankings order differently. Then
\begin{equation}\label{eq:angle}
\E\bigl[d_K(\kappa,\kappa')\bigr]=\binom{e}{2}\frac{\theta}{\pi},\qquad \binom{e}{2}\frac{\theta}{\pi}\le\E\bigl[d_F(\kappa,\kappa')\bigr]\le2\binom{e}{2}\frac{\theta}{\pi}.
\end{equation}

\begin{proof}
Fix coordinates $i<j$ and let $g$ be column $i$ of $W$ minus column $j$, so $g\sim\mathcal N(0,2I)$. The score difference of the pair is then $q^{\top}g$ for $q$ and $q'^{\top}g$ for $q'$. Since $q\neq0$, $q^{\top}g$ is a normal variable with mean zero and positive variance, so the pair ties with probability zero. The same holds for $q'$, so the tie rule plays no part almost surely. The pair is ordered differently exactly when $q^{\top}g$ and $q'^{\top}g$ have opposite signs. If $q'$ is a positive multiple of $q$, the two rankings agree and $\theta=0$. If it is a negative multiple, every pair is reversed and $\theta=\pi$. The formula holds in both cases. Otherwise $q$ and $q'$ span a plane, and we project $g$ onto it. The projection is an isotropic Gaussian vector of that plane, so its direction is uniform on the circle. And only the projection enters $q^{\top}g$ and $q'^{\top}g$. The two signs differ when the line orthogonal to the projection separates $q$ and $q'$. This happens on two arcs of total angle $2\theta$ out of $2\pi$ (Figure~\ref{fig:s-angle}), so with probability $\theta/\pi$ \citep{charikar2002similarity}. Summing over the $\binom{e}{2}$ pairs gives the expected Kendall distance by linearity of expectation, which needs no independence between pairs. Every two rankings satisfy $d_K\le d_F\le2d_K$ \citep{diaconis1977footrule}, and taking expectations gives the bounds.
\end{proof}

% Figure: why a coordinate pair flips with probability theta/pi (supplement, Section S1.6). Drawn by hand; q at angle 0,
% q' at theta = 50 degrees; g flips the pair when it points into (90, 90+theta) or (270, 270+theta) degrees.
\begin{figure}[!htbp]
\centering
\begin{tikzpicture}[>=Stealth, scale=1.15]
\def\th{50}
\draw[black!30] (0,0) circle (2);
% the two arcs of directions of g that separate q and q'
\draw[red!70!black, line width=3pt] (90:2) arc (90:90+\th:2);
\draw[red!70!black, line width=3pt] (270:2) arc (270:270+\th:2);
% q and q'
\draw[->, very thick, blue!70!black] (0,0) -- (0:1.7) node[right] {$q$};
\draw[->, very thick, blue!70!black] (0,0) -- (\th:1.7) node[above right] {$q'$};
\draw (0:0.55) arc (0:\th:0.55);
\node at (\th/2:0.8) {$\theta$};
% one direction of g inside an arc, and the line orthogonal to it
\draw[->, thick, dashed] (0,0) -- (115:1.95) node[above left] {$g$};
\draw[dotted, very thick] (25:2.3) -- (205:2.3);
\node[font=\footnotesize, align=left, anchor=west, red!70!black] at (2.5,1.2) {$g$ in a red arc:\\the pair flips};
\node[font=\footnotesize, align=left, anchor=west] at (2.5,-0.3) {dotted: the line\\orthogonal to $g$};
\end{tikzpicture}
\caption{\textbf{Why a pair of coordinates flips with probability $\theta/\pi$.} The plane spanned by the two feature vectors $q$ and $q'$, which form the angle $\theta$. The projection of the random direction $g$ onto this plane points in a uniformly random direction. The pair is ordered differently for $q$ and $q'$ exactly when the line orthogonal to $g$ (dotted) separates $q$ from $q'$. That happens when $g$ points into one of the two red arcs. The arcs have total angle $2\theta$ out of $2\pi$, so the pair flips with probability $\theta/\pi$. The dashed arrow shows one such $g$.}
\label{fig:s-angle}
\end{figure}


Kendall's rank correlation $\tau$ between two rankings of $e$ items is $1-2d_K/\binom{e}{2}$. So in expectation over $W$, for fixed inputs, $\E[\tau(\kappa,\kappa')]=1-2\theta/\pi$. This is the arcsine kernel of the two feature vectors. It equals twice the arc-cosine kernel of order 0 defined by \citet{cho2009kernel}, minus one. The Kendall kernel of the support vector classifier in Section~\ref{supp:controlled} is this $\tau$, computed for the one projection that each view draws.

The property concerns the random strategy only. The target-aware projection places discriminant coordinates fitted to the labels, and the calibrated projection standardizes the projected coordinates with training statistics. The property holds for each view separately and says nothing about the trained layers above the encoder.

\subsection{Details for Section~\ref*{sec:theory-margin}: Fixed-Layer Margin}\label{supp:margin}

This subsection asks how far an input ranking can move before a fixed ranking layer changes its output. The answer is a margin, much like the score-gap condition of the encoder. A fixed layer with filters $r_1,\dots,r_N$ over $V$ items orders its filters by their responses $D_j=d_F(\pi,r_j)$ to the input $\pi$. Since the footrule is a metric, no response changes by more than the input moves. So, when the responses are distinct, the smallest gap between consecutive responses guarantees the same output for every input in a ball around $\pi$. When responses tie, a second bound at the end of this subsection limits how far the output can move.

\begin{proposition}[Fixed-layer margin]\label{prop:layer-margin}
Let $\pi$ and $\pi'$ be rankings of the $V$ items of a layer with filters $r_1,\dots,r_N$, $N\ge2$.
\begin{enumerate}[label=(\roman*),nosep]
\item For every filter, $|d_F(\pi,r_j)-d_F(\pi',r_j)|\le d_F(\pi,\pi')$.
\item If the responses at $\pi$ are distinct with smallest consecutive gap $g_{\min}$ and $d_F(\pi,\pi')<g_{\min}/2$, the output ordering at $\pi'$ equals that at $\pi$. If only the two smallest responses are separated, by $g_1>0$, then $d_F(\pi,\pi')<g_1/2$ preserves the nearest filter.
\item If the responses at $\pi$ are distinct, let $\Delta^*(\pi)$ be the smallest footrule distance from $\pi$ to an input with a different output ordering. Then $g_{\min}/2\le\Delta^*(\pi)\le D_{(2)}$, with $D_{(2)}$ the second-smallest response.
\end{enumerate}
\end{proposition}

The lower bound is a certified radius, a distance within which the output is guaranteed not to change. The upper bound says that a layer with two distinct filters can always be moved off its output. Under the prototype readout the nearest output filter is the predicted class, so the second part of (ii) is that readout's classification margin. All three parts follow from the reverse triangle inequality of the footrule metric. This inequality says that the distances from two inputs to one filter differ by at most the distance between the inputs.

\begin{proof}[Proof of Proposition~\ref{prop:layer-margin}]
\emph{(i)} This is the reverse triangle inequality for the metric $d_F$,
\[
d_F(\pi,r_j)\le d_F(\pi,\pi')+d_F(\pi',r_j),
\]
together with the same inequality with $\pi$ and $\pi'$ exchanged.

\emph{(ii)} Let $d_F(\pi,\pi')<g_{\min}/2$. For two filters $j$ and $j'$ with $D_j(\pi)<D_{j'}(\pi)$ consecutive in the sorted responses, (i) gives
\[
D_{j'}(\pi')-D_j(\pi')\ge\bigl(D_{j'}(\pi)-d_F(\pi,\pi')\bigr)-\bigl(D_j(\pi)+d_F(\pi,\pi')\bigr)\ge g_{\min}-2\,d_F(\pi,\pi')>0 .
\]
Every consecutive pair keeps its strict order, so the whole ordering is unchanged and no tie arises. For the second claim, let $r_{(1)}$ be the nearest filter at $\pi$ with response $D_{(1)}$ and let $g_1=D_{(2)}-D_{(1)}>0$. If $d_F(\pi,\pi')<g_1/2$, then $D_{(1)}(\pi')<D_{(1)}+g_1/2$. Every other filter $j$ has $D_j(\pi')>D_j(\pi)-g_1/2\ge D_{(2)}-g_1/2=D_{(1)}+g_1/2$. So $r_{(1)}$ remains the unique nearest filter.

\emph{(iii)} By (ii), every input within footrule distance $g_{\min}/2$ of $\pi$ has the same output ordering, so $\Delta^*(\pi)\ge g_{\min}/2$. Distinct responses at $\pi$ force pairwise distinct filters, because equal filters have equal responses at every input. For the upper bound, let $r_{(2)}$ be the filter with the second-smallest response and take $\pi'=r_{(2)}$. Then $d_F(\pi',r_{(2)})=0$. Since the filters are pairwise distinct and $d_F$ is a metric, every other filter has a positive response at $\pi'$. So $r_{(2)}$ is the unique nearest filter at $\pi'$. At $\pi$ the unique nearest filter is $r_{(1)}\neq r_{(2)}$, so the output ordering differs. The set over which $\Delta^*$ is minimized is therefore nonempty. Hence $\Delta^*(\pi)\le d_F(\pi,r_{(2)})=D_{(2)}$.
\end{proof}

This certificate concerns the layer output and the prototype readout. The nearest-neighbor (kNN) readout of Section~\ref{sec:prediction} needs its own condition. Under uniform weights, a positive gap between the $k$th and $(k+1)$th stored distances preserves the whole neighbor set, and hence the vote, for every perturbation below half that gap. Inverse-distance weights need a further bound on the class vote totals and separate treatment of zero distances.

The upper bound needs only pairwise distinct filters. If responses at $\pi$ tie, let $r_{(1)}$ and $r_{(2)}$ be the first two filters of the output ordering at $\pi$, sorted by $(D_j,j)$. Read $D_{(2)}$ as the response of $r_{(2)}$. The argument above then applies unchanged. The output ordering at $\pi'=r_{(2)}$ starts with $r_{(2)}$, and the one at $\pi$ starts with $r_{(1)}$.

Ties limit these certificates. The footrule distance between two rankings is even, because the sum of the positive displacements equals the magnitude of the sum of the negative ones. The largest possible distance, the diameter, is $\lfloor V^2/2\rfloor$ \citep{diaconis1977footrule}. Together with evenness, this leaves at most $\lfloor V^2/4\rfloor+1$ possible values, which is $65$ at $V=16$. A layer with more filters than that has tied responses at every input, by the pigeonhole principle. The principle says that if there are more filters than possible values, some two filters must share a value. Ties do not imply that every tied response changes under a small perturbation. They imply that the certificate of (ii), which needs strict gaps, does not apply to the full ordering.

Even distances also limit the certified radius. Because distances are even, an input other than $\pi$ lies at footrule distance at least $2$ from it. So the full-order certificate of (ii), the one for the whole output ordering, covers an input other than $\pi$ only when $g_{\min}\ge6$. Packing $N$ distinct responses into $[0,\lfloor V^2/2\rfloor]$ with consecutive gaps of at least $6$ also needs $6(N-1)\le\lfloor V^2/2\rfloor$. Among the hidden-layer sizes of ArrowFlow, $V_\ell\in\{16,32,64,128\}$ with $N_\ell\in\{64,128\}$ (Table~\ref{tab:settings}), this fails at $(V_\ell,N_\ell)\in\{(16,64),(16,128),(32,128)\}$, so no input of such a layer carries a full-order certificate. Only $(16,128)$ also has more filters than distance values, so only there does the pigeonhole principle force ties and the hypothesis of (ii) fail outright.

Ties are common well before the pigeonhole bound. For a fixed input and a uniformly random filter, the footrule has mean $(V^2-1)/3$ and variance $(V+1)(2V^2+7)/45$ \citep{diaconis1977footrule}. So its standard deviation grows like $V^{3/2}$, while its range grows like $V^2$. The distances therefore crowd around their mean. At $V=32$ the standard deviation is $38.8$ against a range of $512$. If the $257$ attainable values were equally likely, $128$ random filters would share a distance with another filter with probability $1-(1-1/257)^{127}\approx0.39$. The concentration of the distances raises the share to the $0.788$ that the diagnostics report for uniform random filters (Section~\ref{supp:controlled}).

Pairs of filters tie far less often than responses do. At a fixed input, two independent uniform random filters tie with probability $\rho_V=\sum_k\Pr(D=k)^2$, where $D$ is the footrule distance from the input to a uniform random filter. The distribution of $D$ is close to normal \citep{diaconis1977footrule}, so $\rho_V$ is close to $1/\sqrt{\pi}$ divided by the standard deviation of $D$, and it falls like $V^{-3/2}$. The exact distribution of $D$, which a short recursion over positions gives, yields $\rho_V\approx0.040$, $0.014$, $0.0052$ and $0.0018$ at $V=16$, $32$, $64$ and $128$. It also reproduces the tied shares of uniform random filters in Table~\ref{tab:s-diag-ties}. So at $V=32$ most responses are shared, but only about 1.4 percent of filter pairs tie. With $128$ filters the number of distinct responses is about half the number of filters, so a group of filters at one distance holds about two filters on average.

At $(16,64)$ and $(32,128)$ the responses can be distinct, and Proposition~\ref{prop:layer-margin} applies. But it certifies no input other than $\pi$ itself. ArrowFlow's first hidden layer was one of the three on 89 of the 255 outer folds of the seventeen datasets, 21 at $(16,64)$, 16 at $(16,128)$ and 52 at $(32,128)$, and on 11 of the 17 datasets. At every layer size used, ties are the rule for random filters. A layer of $N$ such filters has on average $\binom N2\rho_V$ tied pairs at an input. This is about $3.7$ at $(V,N)=(128,64)$, the size with the fewest ties, and more than $10$ at each of the other seven sizes. So an input with all responses distinct is rare, and the certificate of (ii) then covers at most the nearest filter. Table~\ref{tab:s-diag-ties} gives the tied shares of the trained layers.

\paragraph{A bound that allows ties.} Part (i) also limits how much of the output ordering can change, whether or not responses tie. Write $\mathrm{out}(\pi)$ for the output ordering at $\pi$, and let $\delta=d_F(\pi,\pi')$. If two filters $j$ and $j'$ are ordered differently in $\mathrm{out}(\pi)$ and $\mathrm{out}(\pi')$, then $|D_j(\pi)-D_{j'}(\pi)|\le2\delta$. Indeed, if $D_j(\pi)<D_{j'}(\pi)-2\delta$, then (i) gives $D_j(\pi')\le D_j(\pi)+\delta<D_{j'}(\pi)-\delta\le D_{j'}(\pi')$, so $j$ comes first at both inputs. The same holds with $j$ and $j'$ exchanged. Hence
\[
d_K\bigl(\mathrm{out}(\pi),\mathrm{out}(\pi')\bigr)\le\#\bigl\{\{j,j'\}:\ |D_j(\pi)-D_{j'}(\pi)|\le2\,d_F(\pi,\pi')\bigr\},
\]
where $d_K$ counts the pairs that two rankings order differently. Since $d_F\le2d_K$ \citep{diaconis1977footrule}, the footrule distance between the two outputs is at most twice this count. Tied responses always enter the count. If the responses are distinct and $d_F(\pi,\pi')<g_{\min}/2$, the count is zero, which is (ii). The output ordering is the input of the next layer or of the readout, so the bound also limits how far that input moves.

\subsection{One Batch Update as a Weighted Borda Count}\label{supp:borda}

This subsection proves Proposition~\ref{prop:batch-borda}, which says what one batch update of a filter computes. The identity prior and the votes give every row of the accumulator the same total mass, which turns the row mean into a weighted position sum. The paragraphs after the proof relate the update to squared position distances and to gradients. They also compare it with the footrule median, the ranking with the smallest weighted sum of footrule distances to the voters.

\begin{proof}[Proof of Proposition~\ref{prop:batch-borda}]
The accumulator $A_j$ has one row per item, in the filter's current order, and one column per position. At the start of the batch it is the identity prior, so the row of item $v$ holds weight one in column $\operatorname{pos}_{r_j}(v)$. A vote of signed weight $a_t$ toward the target $\tilde\tau_t$ adds $|a_t|$ to column $\operatorname{pos}_{\tilde\tau_t}(v)$ of the row of every item $v$~\eqref{eq:vote}. After the batch, the row $p=\operatorname{pos}_{r_j}(v)$ of $v$ has total mass $1+\sum_t|a_t|$, the same for every item. The sum of its column positions, each weighted by the mass in that column, is $\operatorname{pos}_{r_j}(v)+\sum_t|a_t|\operatorname{pos}_{\tilde\tau_t}(v)=P_j(v)$. The mean voted position $s_p$ of~\eqref{eq:rowmean} is the quotient of the two, which is~\eqref{eq:rowmean-borda}. Dividing by a common positive constant does not change the ranking. So the reorder ranks the items by $P_j$ with the fixed tie rule, which is the weighted Borda count~\eqref{eq:borda-update}. The example below shows that the result differs from the footrule median of the same profile.
\end{proof}

\paragraph{The Borda count as a median under squared distance.} The Borda ranking minimizes a weighted sum of squared position differences to the voters. For two rankings $\kappa$ and $\tau$, let $Q(\kappa,\tau)=\sum_v\bigl(\operatorname{pos}_\kappa(v)-\operatorname{pos}_{\tau}(v)\bigr)^2$ be their squared position distance. For voters $\tau_t$ with weights $\omega_t$, let $\sum_t\omega_t Q(\kappa,\tau_t)$ be the weighted squared position distance from $\kappa$ to the voters. When the square is expanded, $\sum_v\operatorname{pos}_\kappa(v)^2$ and $\sum_v\operatorname{pos}_{\tau_t}(v)^2$ do not depend on $\kappa$. So that sum is a constant minus $2\sum_v\operatorname{pos}_\kappa(v)P(v)$, with $P(v)=\sum_t\omega_t\operatorname{pos}_{\tau_t}(v)$. The rearrangement inequality says that a sum of products of two sequences is largest when both are sorted in the same order. So the inner product is largest when positions increase with $P$. Every ranking by $P$, the Borda ranking with the fixed tie rule included, therefore minimizes the weighted squared position distance. The Borda count is thus a median under squared position differences, which makes it a kind of mean. Indeed, the weighted squared distance from $\kappa$ to the voters equals the total weight times the squared distance from the position vector of $\kappa$ to the voters' weighted mean position vector, plus a constant. So the Borda ranking is the ranking whose position vector lies nearest that mean, the mean-proximity reading of the Borda count \citep{zwicker2008average}. The footrule median is its counterpart under absolute differences. The prior enters as one more voter with weight one.

For votes of equal weight, the Borda ranking is also the maximum-likelihood estimate of the center of the distance-based ranking model built on this squared distance \citep{fligner1986distance}. In such a model, a ranking is less likely the farther it lies from the center. The estimate follows because the model's normalizing constant does not depend on the center. The squared position distance measures dissimilarity, but it is not a metric. The Mallows model is a standard model of noisy rankings around a central ranking. Under it, every positional scoring rule with a non-increasing, nonconstant score vector, the Borda count among them, returns the modal ranking with probability tending to one as the number of votes grows \citep{caragiannis2013noisy}. A positional scoring rule awards each position a fixed score and ranks the items by their total scores. The modal ranking is the most likely one. The result assumes a fixed item set, independent votes and a dispersion strictly between zero and one.

\paragraph{The signed objective of one batch.} The same reading extends to negative votes. A negative vote acts as a vote of negative weight for its original target. For a vote of signed weight $a_t$ toward $\tau_t$, $|a_t|\operatorname{pos}_{\operatorname{rev}(\tau_t)}(v)=|a_t|(V+1)+a_t\operatorname{pos}_{\tau_t}(v)$, and the first term is the same for every item. The update of Proposition~\ref{prop:batch-borda} therefore ranks the items by the signed position sum $b(v)=\operatorname{pos}_{r_j}(v)+\sum_t a_t\operatorname{pos}_{\tau_t}(v)$. Centering every position at $(V+1)/2$ gives $Q(\kappa,\tau)+Q(\kappa,\operatorname{rev}\tau)=V(V^2-1)/3$ for every pair of rankings. The update is therefore a minimizer over the $V!$ rankings of $J(\kappa)=Q(\kappa,r_j)+\sum_t a_t Q(\kappa,\tau_t)$. In $J$, a negative vote subtracts squared disagreement with its own target. The set of rankings is finite, so the minimum is attained even with negative weights. When the signed sums tie, the fixed tie rule selects one minimizer among several. The identity and the minimizer property hold in exact arithmetic. The implementation compares computed row means, so it matches them up to rounding error (Section~S2). This objective belongs to one batch with its targets, gates and weights held fixed. It is not an objective of the training run, whose targets change with the network. Nor is it an objective of the readout, which measures footrule neighborhoods.

\paragraph{The returned motion as a gradient.} No derivative is computed anywhere in the ranking layers, but two quantities of the update coincide with gradient quantities. First, the motion that the output vote returns is half a gradient. For an accepted example, let $x$ and $y$ be the position vectors of the last hidden ranking $\pi^{(L-1)}$ and of the filter $o$ of the example's class. Both are indexed by the hidden filters, so that $Q(x,y)=\sum_h(x_h-y_h)^2$ is their squared position distance. The output vote has a nonnegative weight, so its target is $\pi^{(L-1)}$ itself. The motion it returns, keyed by hidden filter, is $\operatorname{pos}_{\pi^{(L-1)}}(h)-\operatorname{pos}_{o}(h)=x_h-y_h$. It is therefore exactly $\tfrac12\nabla_xQ$, half the gradient of the squared position distance with respect to the positions of the hidden filters in the example's ranking. This identity holds by construction. The positions in $x$ are not parameters of the network, and they have no derivative with respect to the hidden filters, which are permutations. A hidden filter with a positive entry would move earlier in $x$ under a gradient step on $Q$. In ArrowFlow it votes toward its input, which tends to move it earlier. So the signal sent to the last hidden layer points the position of each of its filters the way a gradient step on $Q$ would. The vote toward the input thus stands in for the derivative of the sort, as a straight-through estimator does \citep{bengio2013estimating}. Such an estimator trains through a discrete step by replacing the step's derivative with a simple surrogate.

Second, under a condition on the weights, the batch update is one gradient step followed by a projection onto the rankings. Write $x_\kappa$ for the position vector of any ranking $\kappa$ and $L(x)=\tfrac12\sum_ta_t\|x-x_{\tau_t}\|^2$. Whenever $1+\sum_ta_t>0$, the step $x_{r_j}-\eta'\nabla L(x_{r_j})$ with $\eta'=1/(1+\sum_ta_t)$ equals $\eta'b$, with $b$ the signed position sum above. The ranking nearest to that step sorts the items by $b$, so it is the minimizer of $J$ that the update returns. The batch update is then the projection onto the rankings of one gradient step on $L$. On position vectors, this gradient step is a step of learning vector quantization \citep{kohonen1990self}, a classical prototype method, toward attracting targets and away from repelling ones. The projection then maps it to the nearest vertex of the permutahedron, the polytope whose vertices are the position vectors of all rankings. \citet{blondel2020fast} instead project onto the whole polytope, which gives soft ranks that need not be vertices. Neither identity involves differentiation in the implementation, which computes position differences and row means only.

\paragraph{The relayed motion.} A hidden layer $\ell>1$ relays the same kind of quantity to the layer below. Let $x$ be the position vector of its input $\pi^{(\ell-1)}$, and let $y_j$ be that of a selected filter $r_j$ with a nonzero vote $a_j$. Both are indexed by the $V=V_\ell$ items of the layer. The relayed motion $\operatorname{sign}(a_j)\,m_{r_j\to\pi^{(\ell-1)}}$, keyed by item, is $\operatorname{sign}(a_j)(x-y_j)$. So the mean $u$ over the $n$ such filters is exactly half the gradient with respect to $x$ of $\tfrac1n\sum_j\operatorname{sign}(a_j)\,Q(x,y_j)$. This function is the squared distance of the input to the attracting filters, those with positive votes, minus that to the repelling ones. A step against this gradient moves the input toward the attracting filters and away from the repelling ones. The layer below reads $u$ as the last hidden layer reads the output layer's signal. Weighting each term by $|a_j|$ would make $n\,u$ half the gradient of $\sum_ja_jQ(x,y_j)$, the example's own terms in the signed objectives $J$ of the selected filters. The relay weights the filters equally instead.

The earlier relay of Section~\ref{sec:training}, one that did not carry the sign of the vote, sent $m_{r_j\to\operatorname{rev}(\pi^{(\ell-1)})}$ after a repelling vote. Keyed by item, this motion is $(V+1-x)-y_j$. Its inner product with $x-y_j$ is $(V+1)\sum_v(x_v-y_{jv})-\sum_v(x_v^2-y_{jv}^2)=0$ for every pair of rankings, because all rankings of $V$ items share one position sum and one sum of squared positions. To first order, it therefore asked the layer below neither to approach $r_j$ nor to move away from it. Followed alone, it could not move the input away either. Sorting the items by $x-\lambda\bigl((V+1-x)-y_j\bigr)$ with $\lambda>0$ sorts them by $x+\lambda y_j/(1+\lambda)$. This order agrees with $r_j$ on every pair of items on which the input already does.

\paragraph{A sufficient statistic of the accumulator.} The accumulator stores more than the update needs. Every row of $A_j$ carries the same total mass, so only the two row sums enter the reorder. A vector of $V$ signed position scores per filter therefore carries the same information in exact arithmetic, which makes it a sufficient statistic. Storage per layer would fall from $O(NV^2)$ to $O(NV)$. The matrix is kept as it is, for two reasons. First, it also records the full position histogram that a footrule median would need. Second, replacing it would regroup the floating-point sums, which can reverse two nearly tied row means. It would also change the recorded state hashes that seal the completed runs. The observation is algebraic and is not implemented here.

\paragraph{An isolated negative vote.} For the prior and one negative vote alone, the update never decreases the footrule distance to the original target. Raising the vote weight only swaps adjacent items that still agree with the target. For target positions $t_a<t_b$ and adjacent positions $k,k+1$, such a swap cannot lower the footrule distance to the target, since $|k-t_b|+|k+1-t_a|\ge|k-t_a|+|k+1-t_b|$. This says nothing about a sample inside a batch that carries other votes.

\paragraph{Finite-sample disagreement.}
A small profile shows that the Borda count and the footrule median can disagree. Take the prior $ABC$ with weight one and three votes $ABC$, $BCA$, $CBA$ of weight one, so that the profile is $ABC$, $ABC$, $BCA$, $CBA$ with unit weights. The weight condition of Proposition~\ref{prop:borda-iia} holds, $1\le 2\cdot 3$. The position sums are $8$ for $A$, $7$ for $B$ and $9$ for $C$, so the Borda count returns $BAC$. The summed footrule distances from the profile to the six candidate rankings are
\[
\begin{aligned}
&ABC\colon 8,\qquad ACB\colon 12,\qquad BAC\colon 10,\\
&BCA\colon 10,\qquad CAB\colon 14,\qquad CBA\colon 10,
\end{aligned}
\]
so the footrule median is $ABC$, with cost $8$ against $10$ for the Borda ranking. The two rules disagree although the mean positions are distinct. So the disagreement does not come from the tie rule.

The footrule median is an assignment problem, the problem of matching items to positions at the least total cost. For voters $\tau_t$ with weights $\omega_t$, let
\[
\mathrm{cost}(v,q)=\sum_t\omega_t\,\bigl|\operatorname{pos}_{\tau_t}(v)-q\bigr|
\]
be the cost of placing item $v$ at position $q$. A ranking is a perfect matching of items to positions, one item per position, and its cost is the weighted sum of footrule distances to the voters. So a minimum-cost matching computes the median exactly \citep{dwork2001rank}. The prior enters as one more voter with weight one.

\paragraph{The median and the prior.} The footrule median cannot move a prior that outweighs all the votes together. Let the prior $r$ have weight $w_0$ and the votes $\tau_t$ weights $\omega_t$ of total mass $\Omega=\sum_t\omega_t$. Let $F(\kappa)=w_0\,d_F(r,\kappa)+\sum_t\omega_t\,d_F(\tau_t,\kappa)$ be the summed cost of a ranking $\kappa$. The triangle inequality gives $d_F(\tau_t,\kappa)\ge d_F(\tau_t,r)-d_F(r,\kappa)$. Weighting by $\omega_t$ and summing gives $F(\kappa)\ge F(r)+(w_0-\Omega)\,d_F(r,\kappa)$ for every $\kappa$. When $\Omega<w_0$, every $\kappa\neq r$ has $F(\kappa)>F(r)$. So the prior is the unique median, and the median update cannot move the prototype. This is the bound that Section~\ref{supp:laboratory} quotes with its laboratory, an exploratory comparison of aggregation rules.

\paragraph{Disagreement in the limit.}
The two rules can also converge to different rankings. Let independent rankings of $a,b,c$ equal $bac$ with probability $0.6$ and $acb$ with probability $0.4$. The expected positions are $1.6$ for $a$, $1.8$ for $b$ and $2.6$ for $c$, so the ranking by expected position is $abc$. By Proposition~\ref{prop:persistent-estimator}, this is the limit of the persistent Borda estimator, which keeps every vote. The expected footrule distance to a candidate ranking is $1.6$ for $bac$, $2.0$ for $abc$ and larger for the other four. The empirical footrule median converges almost surely to this unique minimizer, $bac$. The reason is that the empirical mean distance to each of the finitely many candidates converges, by the strong law of large numbers. This law says that a sample mean converges to its expected value with probability one. The two rules are therefore consistent for different targets, that is, each converges to its own target as the number of votes grows. The median's limit $bac$ is also the order of the pairwise majorities. A ranking places $b$ before $a$ with probability $0.6$, and it places $a$ and $b$ before $c$ with probabilities $1$ and $0.6$. So the example is an instance of the classical disagreement between the Borda count and the pairwise-majority order of Condorcet. The choice between the Borda count and the footrule median is a design decision (Section~\ref{sec:arrow}).

\subsection{Details for Section~\ref*{sec:theory-aggregation}: The Idealized Persistent Estimator}\label{supp:persistent}

ArrowFlow resets each accumulator after every batch, so it keeps no vote from an earlier batch. This subsection asks what an idealized estimator that kept every vote would converge to. Section~\ref{sec:theory-aggregation} states in words that it would converge to the ranking by expected position. This is the statement.

\begin{proposition}[Idealized persistent estimator]\label{prop:persistent-estimator}
Let $\pi_1,\pi_2,\dots$ be independent rankings of $V$ items drawn from one distribution with pairwise distinct expected positions $\mu_v=\E[\operatorname{pos}_{\pi_t}(v)]$. For a fixed prior ranking $r$ and prior weight $w_0\ge0$, let $\hat\mu_T(v)=\bigl(w_0\operatorname{pos}_r(v)+\sum_{t=1}^{T}\operatorname{pos}_{\pi_t}(v)\bigr)/(w_0+T)$. Then, almost surely, the ranking by $\hat\mu_T$ equals the ranking by $\mu$ for all large $T$.
\end{proposition}

The proof applies the strong law of large numbers to the position of each item.

\begin{proof}[Proof of Proposition~\ref{prop:persistent-estimator}]
By the strong law of large numbers, $T^{-1}\sum_{t=1}^{T}\operatorname{pos}_{\pi_t}(v)\to\mu_v$ almost surely for every item $v$. The prior term $w_0\operatorname{pos}_r(v)/(w_0+T)$ tends to zero and $T/(w_0+T)\to1$, so $\hat\mu_T(v)=\frac{w_0\operatorname{pos}_r(v)}{w_0+T}+\frac{T}{w_0+T}\cdot\frac1T\sum_{t=1}^{T}\operatorname{pos}_{\pi_t}(v)\to\mu_v$ almost surely. There are finitely many items, so the convergences hold simultaneously. For $V=1$ the claim is immediate. For $V\ge2$ let $\gamma>0$ be the smallest difference between two expected positions. Almost surely there is a $T_0$ such that $|\hat\mu_T(v)-\mu_v|<\gamma/3$ for every $v$ and every $T\ge T_0$. For any two items with $\mu_u<\mu_v$ this gives $\hat\mu_T(v)-\hat\mu_T(u)>\gamma-2\gamma/3>0$. So every pair keeps the order of its expectations, and no tie occurs.
\end{proof}

The estimator keeps every vote with equal weight and draws the votes from one fixed distribution. The network differs in three ways. First, it resets its accumulator to the identity prior after every batch. Second, it weights and accepts votes according to the current predictions. Third, it changes the representations that produce the votes. The proposition therefore characterizes the aggregation rule under fixed conditions. It is not an end-to-end convergence result for the training procedure.

\paragraph{The reset batch update.} The reset changes the outcome even for votes from one fixed distribution. Draw batches of $16$ votes of weight $1/20$ from the distribution of the limit example above, which gives $bac$ with probability $0.6$ and $acb$ with probability $0.4$. With the prior of weight one, these batches meet the weight condition of Proposition~\ref{prop:borda-iia}. If the filter is $bca$ and $k$ of the votes are $bac$, then $P(c)-P(b)=(4+3k)/20>0$ and $P(a)-P(c)=1/5>0$. So the update returns $bca$ for every batch, although the ranking by expected position is $abc$. Each update depends only on the current filter and a fresh batch. So the filters form a Markov chain, a random process whose next state depends only on its current state. The exact Markov chain of this update has three absorbing rankings, $abc$, $bca$ and $cab$, which it never leaves once it enters them. So its limit depends on the initial filter.
